%\documentclass[twocolumn]{article}\
\documentclass{article}
\usepackage{amsmath, amssymb, cancel, mathtools, bm}
\usepackage[left=.5in, right=.5in, top=1in, bottom=1in]{geometry}
\usepackage[most]{tcolorbox}
\usepackage[skip=1em,indent=0pt]{parskip}
\setlength{\parindent}{0pt}
\newcommand{\statvec}[1]{\underset{\sim}{\bm{#1}}} % Vector symbol (tilde under vector)
\DeclareMathOperator{\EX}{\mathbb{E}} % Expected Value symbol
\DeclareMathOperator{\ext}{\EX_{\theta}} % Expected Value symbol
\DeclareMathOperator{\vart}{Var_{\theta}} % Expected Value symbol
\newcommand{\indep}{\perp\!\!\!\!\perp} % Independence symbol
\newcommand{\real}{\mathbb{R}} % Simplified 'Reals' indicator
\newcommand{\pto}{\overset{P}{\to}}
\newcommand{\asto}{\overset{a.s.}{\to}}
\newcommand{\rthto}{\overset{L^r}{\to}}
\newcommand{\dto}{\overset{D}{\to}}
\newcommand{\xtb}{x^{\top}\statvec{\beta}}
\newcommand{\xitb}{x_i^{\top}\statvec{\beta}}
% Force all aggregate symbols to always put values above/below
\let\Oldint=\int
\let\Oldsum=\sum
\let\Oldprod=\prod
\let\Oldbigcup=\bigcup
\let\Oldbigcap=\bigcap
\let\Oldlim=\lim
\renewcommand{\int}{\Oldint\limits} 
\renewcommand{\sum}{\Oldsum\limits} 
\renewcommand{\prod}{\Oldprod\limits} 
\renewcommand{\bigcup}{\Oldbigcup\limits} 
\renewcommand{\bigcap}{\Oldbigcap\limits} 
\renewcommand{\lim}{\Oldlim\limits} 
\newcommand{\infint}{\int_{-\infty}^{\infty}}
\newcommand{\iiddist}{\overset{\mathrm{iid}}{\sim}}
\newcommand{\norm}[1]{\left\lVert#1\right\rVert}
\newcommand{\boldred}[1]{\textbf{\textcolor{red}{#1}}}


\begin{document}
%%%%%%%%%%%%%%%%%%%%%%%%%%%%%%%%%%%%%%%%%%%%%%%%%%
\section{}

$ A = \begin{pmatrix} 1 & 3 \\ 5 & 2 \end{pmatrix}, B = \begin{pmatrix} 4 & 9 \\ 2 & 6 \end{pmatrix}$

a. $A + B = \begin{pmatrix} 1+4 & 3+9 \\ 5+2 & 2+6 \end{pmatrix} = \begin{pmatrix} 5 & 12 \\ 7 & 8 \end{pmatrix}$

b. $A - B = \begin{pmatrix} 1-4 & 3-9 \\ 5-2 & 2-6 \end{pmatrix} = \begin{pmatrix} -3 & -6 \\ 3 & -4 \end{pmatrix}$

c. $AB = \begin{pmatrix} 1(4)+3(2)& 1(9)+3(6) \\ 5(4)+2(2) & 5(9)+2(6) \end{pmatrix} = \begin{pmatrix} 10 & 27 \\ 24 & 57 \end{pmatrix}$

d. $BA = \begin{pmatrix} 4(1)+9(5) & 4(3)+9(2) \\ 2(1)+6(5) & 2(3)+6(2) \end{pmatrix} = \begin{pmatrix} 49 & 30 \\ 32 & 18 \end{pmatrix}$

%%%%%%%%%%%%%%%%%%%%%%%%%%%%%%%%%%%%%%%%%%%%%%%%%%
\section{}
$ A = \begin{pmatrix} 9 & 1 \\ 1 & 4 \end{pmatrix}, B = \begin{pmatrix} 2 & 6 & 5 \\ 3 & 1 & 8 \end{pmatrix}$

a. Non-conformable

b. Non-conformable

c. $AB = \begin{pmatrix} 9(2)+1(3) & 9(6)+1(1) & 9(5)+1(8) \\ 1(2)+4(3) & 1(6)+4(1) & 1(5)+4(8) \end{pmatrix} = \begin{pmatrix} 21 & 55 & 53 \\ 14 & 10 & 37 \end{pmatrix}$

d. Non-conformable

%%%%%%%%%%%%%%%%%%%%%%%%%%%%%%%%%%%%%%%%%%%%%%%%%%
\section{}
Show that the following is true for any invertible, square matrics A,B,C,D of the same size: $[(A+B)^{\top}D^{-1}C^{-1}]^{-1} = CD(A^{\top}+B^{\top})^{-1}$

By the distribution of inverses (applying it twice), then distributing the transpose, we get:
$[(A+B)^{\top}D^{-1}C^{-1}]^{-1} = (D^{-1}C^{-1})^{-1}[(A+B)^{\top}]^{-1} = CD[(A+B)^{\top}]^{-1} = CD(A^{\top} + B^{\top})^{-1}$

%%%%%%%%%%%%%%%%%%%%%%%%%%%%%%%%%%%%%%%%%%%%%%%%%%
\section{}
$A = \begin{pmatrix} A_{11} & A_{12} \\ A_{21} & A_{22} \end{pmatrix}, B = \begin{pmatrix} B_{11} & B_{12} \\ B_{21} & B_{22} \end{pmatrix}$

a. $AB = \begin{pmatrix} A_{11}B_{11}+A_{12}B_{21} & A_{11}B_{12}+A_{12}B_{22} \\ A_{21}B_{11}+A_{22}B_{21} & A_{21}B_{12}+A_{22}B_{22} \end{pmatrix}$

b. To tranpose a partioned matrix, you transpose each partion in the overall matrix, then each partion gets transposed in itself

$A^{\top} = \begin{pmatrix} A_{11}^{\top} & A_{21}^{\top} \\ A_{12}^{\top} & A_{22}^{\top} \end{pmatrix}$

%%%%%%%%%%%%%%%%%%%%%%%%%%%%%%%%%%%%%%%%%%%%%%%%%%
\section{}
$x = \begin{pmatrix} x_1 \\ x_2 \end{pmatrix}, A = \begin{pmatrix} a_{11} & a_{12} \\ a_{12} & a_{22} \end{pmatrix}$

Find $\frac{d q(x)}{dx}$ where $q(x) = x^{\top}Ax$

Treating this more like normal algebra, $x^{\top}Ax$ would be the equivalent to $Ax^2$, and $\frac{d}{dx} Ax^2 = 2Ax$

Thus, because $A$ is symmetric: 
$\frac{d q(x)}{dx} = 2Ax = \begin{pmatrix} a_{11}x_1+a_{12}x_2 \\ a_{12}x_1+a_{22}x_2 \end{pmatrix}$ 

\end{document}